\section{Critical constants for rectangular mean observations}
For a disjoint partition $\mathcal P$ of a finite counting space, its
normalized means are $(Cf)_R=|R|^{-1/2}\sum_{x\in R}f(x)$. They satisfy
$CC^*=\id$ and $C^*C=P$, the orthogonal projection onto functions constant
on each cell. This concrete normalization matters for the penalty.

\begin{theorem}\label{thm:critical}
Let $L=L^*$ on a finite real Hilbert space, $c>0$, and $C$ a coisometry
onto $\R^q$, $q\geq1$. On $\mathcal K=\ker C$ and
$\mathcal O=\im C^*$, write
\[
 L-c\id=\begin{pmatrix}A_0&X\\X^*&B_0\end{pmatrix}.
\]
A finite $k\geq0$ satisfying $L+kC^*C\geq c\id$ exists if and only if
\begin{equation}\label{eq:compatibility}
 A_0\geq0,\qquad X^*(\ker A_0)=0.
\end{equation}
When these hold, its smallest value is
\[
 k_{\min}=\max\{0,\lambda_{\max}(X^*A_0^\dagger X-B_0)\}.
\]
Here $A_0^\dagger$ is the Moore--Penrose inverse on $\mathcal K$.
\end{theorem}
\begin{proof}
The restriction to $\mathcal K$ forces $A_0\geq0$. For $z\in\ker A_0$
the augmented form at $tz+y$ has linear term $2t\langle z,Xy\rangle$
and no quadratic term in $t$. Nonnegativity for every real $t$ forces
$X^*z=0$. This is equivalent to $\im X\subseteq\im A_0$. Under that
condition the full form at $z+y$ is exactly
\[
 \langle z+A_0^\dagger Xy,A_0(z+A_0^\dagger Xy)\rangle
 +\langle y,(B_0+k\id-X^*A_0^\dagger X)y\rangle.
\]
Finite spectral decomposition proves the criterion and the minimum.
\end{proof}
This is the classical generalized Schur criterion~\cite[Theorem 1]{Higham},
applied to actual means; it is not a new general matrix theorem. Strict
positivity of $A_0$ suffices, whereas a weak bound at the target requires
the extra compatibility in~\eqref{eq:compatibility}.

\begin{example}\label{ex:rectangle}
On $(\mathbb Z/2\mathbb Z)^3$ take the same scalar coefficient in every
gradient direction,
\[
 a(x)=6+(-1)^{x_1}+4(-1)^{x_2},\qquad L=d^*ad.
\]
Its values are $1,3,9,11$, each twice, and its mean is $m=6$; hence
$1\leq a\leq2m$. Choose the four actual rectangles
$R_{yz}=\{0,1\}\times\{y\}\times\{z\}$ and
$(Cf)_{yz}=(f(0,y,z)+f(1,y,z))/\sqrt2$.
The best coercivity constant of $L$ on $\ker C$ is $16$, but no finite
quadratic penalty for these means attains the same constant on the full
space. Every $0<c<16$ is attainable. Thus the supremum $16$ is not attained.
\end{example}
\begin{proof}
For $S\subseteq\{1,2,3\}$ let
$\chi_S(x)=(-1)^{\sum_{i\in S}x_i}/\sqrt8$. These are an orthonormal real
Walsh basis. Retain all $24$ oriented forward rows, including opposite rows
at period two. The actual adjoint calculation gives
\begin{align*}
 L\chi_S={}&24|S|\chi_S
 +4\bigl(|S|-\boldsymbol1_{1\in S}\bigr)\chi_{S\triangle\{1\}}\\
 &+16\bigl(|S|-\boldsymbol1_{2\in S}\bigr)\chi_{S\triangle\{2\}}.
\end{align*}
Indeed $d_i\chi_S=-2\chi_S$ when $i\in S$, and the corresponding
adjoint term is $2(a+T_i a)\chi_S$; absent directions contribute zero.
In the kernel basis $(\chi_1,\chi_{12},\chi_{13},\chi_{123})$ the compression
has the two blocks
\[
 \begin{pmatrix}24&16\\16&48\end{pmatrix},\qquad
 \begin{pmatrix}48&32\\32&72\end{pmatrix}.
\]
After subtracting $16\id$, the first is positive semidefinite with kernel
vector $(-2,1)$, and the second is positive definite, with first pivot $32$
and determinant $768$. This proves the complete kernel bound, not just a
bound on selected test vectors.

Set $v=-2\chi_1+\chi_{12}$ and $u=\chi_2$. Direct evaluation gives
\[
 Cv=0,\quad\|v\|^2=5,\quad\|Cu\|^2=\|u\|^2=1,
 \quad(L-16\id)v=4u,\quad\langle u,(L-16\id)u\rangle=8.
\]
The augmented form on $u+tv$ is therefore $8+8t+k$, negative for
$t<-(8+k)/8$. No finite $k$ works. For $c<16$ the compressed block is
strictly positive and Theorem~\ref{thm:critical} supplies a finite penalty;
for $c>16$, $v$ violates coercivity already in the kernel.
\end{proof}

\paragraph{A quantitative penalty obstruction.}
Write $\delta=16-c>0$. On the orthonormal pair $(v/\sqrt5,u)$ the
augmented form has matrix
\[
 \begin{pmatrix}\delta&4/\sqrt5\\4/\sqrt5&8+\delta+k\end{pmatrix}.
\]
Its determinant gives the necessary bound
\begin{equation}\label{eq:blowup}
 k\geq\max\{0,16/(5\delta)-8-\delta\}.
\end{equation}
This lower bound need not equal the full optimum. It proves divergence
as the target approaches $16$ from below without numerical extrapolation.
If $F(0)=0$, $DF(0)=C$ and $0\leq\Phi(y)\leq k\|y\|^2$ near zero,
local coercivity would imply $L+kC^*C\geq c\id$ by substituting $sf$,
dividing by $s^2$ and taking $s\to0$. Higher order terms with this same
derivative do not repair the failure at $c=16$.

\section{Transfer to actual flat coexact cochains}
The preceding obstruction also holds on the actual finite coexact domain,
with a proved form dictionary. Fix $N_c\geq2$ and an actual coordinate unit
$e\in\R^{N_c^2-1}$. At the trivial background, a physical one-cochain is
$A_j(x)\in\R^{N_c^2-1}$, $j=1,2,3$, with counting inner product. Use one
unordered plaquette copy for each $i<j$:
\[
 (d_1 A)_{ij}=d_iA_j-d_jA_i,\qquad
 \operatorname{div}A=\sum_j(A_j-T_j^{-1}A_j).
\]
Expansion of curl squares and cyclic summation by parts cancel the cross
terms against divergence, proving the all-cochain identity
\[
 \|d_1A\|^2+\|\operatorname{div}A\|^2=\sum_{i,j,x}\|d_iA_j(x)\|^2.
\]
These are the actual repository conventions, not a supplied comparison.
Take the actual componentwise repair coefficient $w(i,(x,j))=a(x)-1\geq0$.
The total Hodge-plus-repair form is consequently
\begin{equation}\label{eq:physical-dictionary}
 Q_w(A)=\|d_1A\|^2+\|\operatorname{div}A\|^2+
 \sum_{i,j,x}w(x)\|d_iA_j(x)\|^2
 =\sum_{i,j,x}a(x)\|d_iA_j(x)\|^2.
\end{equation}
It is a direct sum of $3(N_c^2-1)$ scalar copies. Componentwise normalized
means on the same four bars have output dimension $12(N_c^2-1)$, and
the whole ambient kernel has coercivity constant $16$ by the two blocks
in Example~\ref{ex:rectangle}. Restriction to coexact fields preserves the
inequality; coisometry on the restricted domain is not assumed.

\begin{proposition}\label{prop:coexact}
At target $16$, no finite quadratic penalty for these actual componentwise
rectangle means supplies coercivity even on $\im d_1^*$. The same is true
for local nonlinear observations with derivative equal to those restricted
means and a locally quadratic penalty. On that domain every $0<c<16$ is
attainable, and the necessary bound~\eqref{eq:blowup} still applies.
\end{proposition}
\begin{proof}
Construct actual fields $U_3=ue$, $V_3=ve$ from Example~\ref{ex:rectangle},
with all other components zero. They are isometric copies of $u,v$ and
independent of $x_3$, so their actual backward divergence vanishes.
More explicitly, use unordered two-cochains with only
\[
 B^U_{23}=-\tfrac12\chi_2 e,\qquad
 B^V_{13}=(\chi_1-\tfrac12\chi_{12})e
\]
nonzero. At side two $d_i^*=T_i^{-1}-\id=T_i-\id$. The adjoint of the
displayed curl puts $d_i^*B_{ij}$ in component $j$ and $-d_j^*B_{ij}$ in
component $i$. Since $d_3^*B^U=d_3^*B^V=0$ and each relevant character
has multiplier $-2$, these actual potentials satisfy $d_1^*B^U=U$ and
$d_1^*B^V=V$. Thus the entire line $U+tV$ is coexact.

Equation~\eqref{eq:physical-dictionary} and the actual means give the same
augmented value $8+8t+k$. Scaling a violating coexact field towards zero
also excludes locally bounded quadratic penalties and the stated nonlinear
observations. Strict positivity on the ambient kernel supplies a finite
penalty for every $c<16$, which can then be restricted. The coexact field
$V$ excludes $c>16$. The same two-dimensional determinant proves the bound.
\end{proof}

\paragraph{Scope and source of the transfer.}
The actual form identity above agrees with
\texttt{componentwiseFlatGradient\_eq\_hodge} and
\texttt{actualTotalStiffness\_energy} in the pinned repository, respectively
in \path{YangMills/RG/ComponentwiseFlatHodge.lean} and
\path{YangMills/RG/ActualPacketSupportEnvelope.lean}. Their proofs and the
actual curl/divergence definitions were read; no new Lean compilation is
claimed. Unit 166 checks the newly constructed physical stencils rather
than relabelling a scalar diagnostic. For derivatives $nd_i$ at $n=2$,
energies and targets scale by four, so $16$ becomes $64$ with rescaled
potentials. This is finite rescaling, not a physical mass-gap claim.

The example has finite inverse moments at this one size. It refutes weak
kernel sufficiency for these prescribed rectangles at the same target;
it does not rule out another rectangular partition with the desired rank,
or establish an unbounded-size violation of the critical rank-order
conjecture. Only the trivial-background form has been identified here.

